\midsectiontitle{Combate Corpo a Corpo}

No campo de batalha de \textit{Overgrown}, a supremacia do combate corpo a corpo pertence àqueles que dominam não apenas a lâmina, mas a própria fúria que ela carrega. As \textbf{Habilidades de Combate} são técnicas ativas e instintivas que podem ser acionadas durante o combate gastando \textbf{Pontos de Combate (PC)}, a reserva de agressividade e foco físico de um guerreiro.

\begin{rpgboxtips}{Pontos de Combate}
Os PC se recuperam apenas com descanso longo. Algumas habilidades na \textbf{Árvore de Talentos} aumentam seu máximo ou permitem recuperação parcial durante o combate. Saber quando gastar e quando poupar é tão importante quanto saber atacar.
\end{rpgboxtips}

\section*{Habilidades}

\combatSkill{Fúria}{2}{Você entra em um estado de \textit{Luta ou Fuga} instintivo. Seu personagem fica compelido a atacar quem considera inimigo, ganhando \textbf{mais um ataque por turno} e \textbf{RD contra danos físicos} durante \textbf{1D6+1 rodadas}. Toda lógica cede ao instinto.}

\combatSkill{Fúria Descontrolada}{3}{Requer \textbf{Fúria ativa}. Ascende o estado de fúria ao limite absoluto: aumenta a duração em \textbf{+3 rodadas}, concede \textbf{mais um ataque por turno} e estende a \textbf{RD para danos mágicos}. A violência desse estado cobra um custo em vida a cada rodada que permanece ativo.}

\combatSkill{Pontos Fatais}{3}{Gaste sua \textbf{Reação} para realizar um teste de \link{per:discernimento}{Discernimento} contra percepção do alvo. Ao passar, você identifica uma abertura em suas defesas, reduzindo sua \textbf{RD apenas para os seus ataques}. Se a redução ultrapassar a RD total, você também ganha \textbf{+2 na Margem de Ameaça} contra esse alvo. Por padrão passar com até 10, tira 1 de RD, passar com 11-15 tira 3 de RD, passar com 16-20, }

\combatSkill{Hemorragia Interna}{1}{Ao acertar um ataque em um inimigo que já está \link{status:sangramento}{Sangrando}, você aprofunda o ferimento girando a lâmina. O alvo sofre imediatamente \textbf{1 tique do sangramento como dano extra}, adicional ao dano do ataque.}

\combatSkill{Garrote}{2}{O primeiro \textbf{Ataque Furtivo} bem-sucedido contra um alvo o deixa \textbf{Silenciado} por \textbf{2 turnos}, incapaz de conjurar magias ou emitir sons que pudessem alertar aliados.}

\combatSkill{Emboscar}{3}{O primeiro \textbf{Ataque Furtivo} bem-sucedido contra um alvo inflige \textbf{1,5× o dano normal}. A surpresa transforma o primeiro golpe em uma execução quase perfeita.}

\combatSkill{Marcar para a Morte}{4}{Você foca toda sua letalidade em um único alvo. Ganha bônus em \textbf{acerto}, \textbf{Margem de Ameaça} e \textbf{dano} contra ele. Em contrapartida, o dano recebido de todas as outras fontes aumenta e você não pode \textbf{desviar ou resistir} ao que não vem do alvo marcado. \par\textit{\footnotesize \textbf{Aprimorado:} As penalidades são progressivamente removidas conforme a habilidade é refinada na Árvore de Talentos.}}

\combatSkill{Última Resistência}{5 por turno}{Ao receber um dano que seria fatal, ao invés de entrar em \textbf{Morrendo}, você permanece com \textbf{1 de vida}, zerando seus PC. Role \textbf{1D4+X}: durante esses turnos você está \textbf{imune à condição Morrendo}. Sobreviver tem um preço.}

\combatSkill{Vício em Sangue}{2}{Se um alvo com \textbf{metade da vida ou menos} estiver em alcance médio, você se arremessa até ele \textbf{sem gastar sua Ação de Movimento}. O cheiro de sangue aguça o instinto predatório.}

\combatSkill{Sede de Sangue}{3}{Ao \textbf{executar} um inimigo, você canaliza o ímpeto da matança em vitalidade. Cure uma quantidade de pontos de vida igual ao \textbf{dano do seu último ataque} contra o alvo abatido.}

\combatSkill{Instinto Sanguinário}{4}{Você abandona a capacidade de \textbf{Reagir} em troca de instinto puro. Ganha bônus de dano progressivo baseado no estado físico de \textbf{todos} os envolvidos no combate — quanto mais feridos todos estiverem, mais letal você se torna.}

\combatSkill{Elipse Carmesim}{3 / 5}{Você gira e corta tudo ao seu redor: causa o \textbf{dano da arma em todos os inimigos} em alcance curto que estejam \link{status:sangramento}{Sangrando}. \par\textit{\footnotesize \textbf{Aprimorado (5 PC):} O ataque pode resultar em \textbf{crítico em cada inimigo acertado}, não apenas no alvo principal.}}

\combatSkill{Ataque Secundário}{2}{Gaste PC para realizar \textbf{um segundo ataque} contra o mesmo alvo no mesmo turno, como uma ação adicional dentro da sua Ação Padrão.}

\combatSkill{Desarranjar}{2}{Realiza uma manobra de afastamento que nega \textbf{ataques de oportunidade} ao se reposicionar, quebrando o engajamento sem expor a guarda.}

\combatSkill{Pisada de Empacto}{3}{(Precisa está em Furia ativa) Gaste PC para realizar um ataque de área que atinge o chão, causando uma onda de choque. Causa 1D6+MIGHT.}

\section*{Mecânicas de Esforço}

As Mecânicas de Esforço são manobras e técnicas táticas disponíveis a qualquer combatente. Diferente das Habilidades, não exigem especialização, apenas PC e o momento certo.

\begin{rpgboxtips}{Quando Usar}
As mecânicas de esforço são opções abertas para todos os personagens, independente de estilo. Combiná-las com as habilidades certas pode transformar uma troca de golpes comum em uma vitória dominante.
\end{rpgboxtips}

\combatSkill{Agarrar}{2}{Realiza uma manobra de imobilização. Teste de \textbf{Força} contra o teste de \textbf{Esquiva} do alvo. Alvos agarrados \textbf{não podem se mover}, sofrem \textbf{Desvantagem} em ataques e precisam passar num teste de Força ou ser auxiliados por aliados para se libertar.}

\combatSkill{Impulsão}{1}{Gaste PC para ganhar \textbf{+5 metros} na sua Ação de Movimento neste turno, canalizando um burst de energia física para cobrir mais terreno.}

\combatSkill{Contra-Ataque}{2}{Se um inimigo \textbf{errar um ataque} contra você, use sua \textbf{Reação} para responder imediatamente com um ataque contra ele.}

\combatSkill{Recuar Estratégico}{2}{Ao desengajar, use esta técnica para se afastar o \textbf{dobro da distância} normal, criando uma zona de segurança entre você e o oponente.}

\combatSkill{Inabalável}{2}{Use como \textbf{Reação} para consumir PC e prevenir as condições \link{status:atordoado}{Atordoado} ou \link{status:derrubado}{Caído} antes que entrem em efeito.}

\combatSkill{Estocada Interruptora}{3}{Use sua \textbf{Reação} para interromper uma habilidade ou magia em execução pelo inimigo, forçando-o a gastar a ação sem efeito.}

\combatSkill{Comando de Assalto}{2}{Gaste sua \textbf{Ação de Movimento} para emitir uma ordem tática a um aliado. Ele realiza \textbf{um ataque fora do seu turno} contra o inimigo de sua escolha.}

\combatSkill{Analisar Padrão}{3}{Gaste seu \textbf{turno inteiro} estudando os movimentos do inimigo. Até o final do próximo turno, todos os ataques contra ele têm \textbf{Vantagem}.}

\combatSkill{Postura de Muralha}{1 por rodada}{Você não pode atacar, mas \textbf{dobra sua Resistência} até o início do seu próximo turno. Uma defesa impenetrável para quem sabe que o maior dano é o recebido.}

\combatSkill{Desarmar}{2}{Tente arrancar a arma do alvo. Teste de \textbf{Força} contra o teste de \textbf{Força} do alvo. Se bem-sucedido, a arma cai no terreno adjacente.}

\combatSkill{Rasteira}{2}{Tente derrubar o alvo aplicando a condição \link{status:derrubado}{Derrubado}. Teste de \textbf{Esquiva} do alvo contra o seu \textbf{teste de Combate}.}

 

